\section{Related Work}
\label{sec:related}

\para{Linear representation hypotheses.} LRH work asks when model representations admit approximately linear concept structure. Park \etal formalize LRH with counterfactual pairs, concept directions, and a causal inner product that better matches semantic structure than naive Euclidean comparisons \cite{park2024lrh}. Follow-up work extends this geometry to categorical and hierarchical concepts, showing that richer concept families can still occupy organized linear structure \cite{park2025categorical}. These papers establish the clean, separable regime that motivates our control experiment. Unlike them, we focus on whether the {\it path} between endpoints remains meaningful when those endpoints are not cleanly separable.

\para{Ambiguous and non-binary concepts.} The assumption of a single clean counterfactual direction becomes harder to maintain for ambiguous or entangled concepts. Nguyen and Leng propose a more flexible maximum-likelihood framework for LRH-style estimation, arguing that realistic concept variation often requires more than a strict paired-difference estimator \cite{nguyen2025sand}. Our findings are compatible with that view: the failure mode we observe is not binary contradiction in the abstract, but changes that disrupt relational structure while keeping surface form nearby.

\para{Interpolation geometry and manifold structure.} A separate line of work studies when straight-line interpolation is a poor proxy for geodesic traversal. Shao \etal derive the pull-back geometry of deep generative models and show how to compute curves that respect the learned manifold \cite{shao2017riemannian}. Struski \etal generalize interpolation to feature-aware geodesics under arbitrary latent densities \cite{struski2019feature}. Michelis and Becker show that linear interpolation can be arbitrarily suboptimal even when endpoints look reasonable \cite{michelis2021linear}. This literature motivates our path metrics, but it says little about text representations or semantic incongruence. We bring these geometric ideas into a text benchmark with matched congruent and incongruent endpoints.

\para{Positioning.} Our work bridges these themes. We inherit the control logic of LRH, but our central object is the interpolation trajectory rather than the direction estimate. We borrow manifold-sensitive evaluation from latent interpolation research, but we apply it to text embeddings and analyze semantic mismatch types directly. The main empirical payoff is a more precise claim: not all incongruence destabilizes interpolation. Relation changes do; short lexical swaps often do not.
